\documentclass[11pt]{article}
\usepackage{fontspec}
\usepackage[french]{babel}
\usepackage{amsmath}
\usepackage{amssymb}
\usepackage{amsthm}
\usepackage{mathrsfs}
\usepackage{tikz-cd}
\usepackage{enumitem}

\newcommand{\printedpage}[1]{\marginpar{\footnotesize p.~#1}}
\newcommand{\msp}[1]{\marginpar{\footnotesize [#1]}}
\newcommand{\Esp}{\operatorname{Esp}}
\newcommand{\Fl}{\operatorname{Fl}}
\newcommand{\Ob}{\operatorname{Ob}}
\newcommand{\Hom}{\operatorname{Hom}}
\newcommand{\Ho}{\operatorname{Ho}}
\newcommand{\id}{\mathrm{id}}
\newcommand{\pr}{\mathrm{pr}}
\newcommand{\M}{\mathcal{M}}

\begin{document}

\printedpage{29}
% head: A. Grothendieck : Dérivateurs, Ch. XVII, Édition des Universités de Montpellier II et Paris VII
% foot: Version 9 février 2019
lesquels sont visiblement adjoint l'un de l'autre (puisque induits par un couple de foncteurs % [sic] adjoint -> adjoints, accord avec « lesquels »
adjoints $u_!$, $u^*$), et de plus compatibles aux localiseurs. Donc on peut appliquer le lemme des foncteurs adjoints pour les localisés. Pour résumer :

\textbf{Corollaire 3.5.} \textit{Supposons que $I$, $I'$ soient artiniens, et qu'on soit dans l'un des deux cas suivants :}
\begin{enumerate}
\item[a)] \textit{Les espaces} $\Esp(I)$, $\Esp(I')$ \textit{sont noethériens (p.ex. $I$, $I'$ finis).}
\msp{page 58}
\item[b)] \textit{$C$ et $TC = C\cap W$ sont stables par composition transfinie.}
\end{enumerate}
\textit{Supposons de plus que $u^*$ applique $\M(I'^\circ)_{\mathrm{c}}$ dans $\M(I^\circ)_{\mathrm{c}}$, ce qui est le cas notamment si $u$ est localement cofiltrant, i.e. les ensembles ordonnés localisés $i'\backslash I$ sont localement cofiltrants. Alors le foncteur $u_!^W$ construit dans (3.3.1) est adjoint à gauche de $u_W^*$.}

\textbf{Remarque 3.6.} Quand on a un couple de foncteurs adjoints
\[
\begin{tikzcd}
\mathcal{C} \arrow[r, shift left, "u_!"] & \mathcal{C}' \arrow[l, shift left, "u^*"]
\end{tikzcd}
\ ,
\]
avec les localiseurs $W$, $W'$ de $\mathcal{C}$, $\mathcal{C}'$, $u^*$ compatible aux localiseurs, et si de plus on a une sous-catégorie pleine $\mathcal{C}_0$ de $\mathcal{C}$ telle que
\begin{enumerate}
\item[a)] $u_!|\mathcal{C}_0$ transforme flèches de $W$ en flèches de $W'$, et
\item[b)] l'inclusion $\mathcal{C}_0 \hookrightarrow \mathcal{C}$ induit une \textit{équivalence} $\mathcal{C}_0W_0^{-1} \xrightarrow{\sim} \mathcal{C}W^{-1}$ (où $W_0 = W\cap\Fl(\mathcal{C}_0)$),
\end{enumerate}
on trouve deux foncteurs canoniques
\[
\begin{tikzcd}
\mathcal{C}W^{-1} \arrow[r, shift left, "u_!^W"] & \mathcal{C}'W'^{-1} \arrow[l, shift left, "u_W^*"]
\end{tikzcd}
\ ,
\]
où $u_!^W$ est défini seulement à isomorphisme canonique près.

\msp{page 59}
Contrairement à ce que prétend \textsc{Anderson}, il n'est par vrai en général que ces deux % [sic] par -> pas, faute de frappe
foncteurs soient adjoints, sans hypothèses supplémentaires. Voici un contrexemple : on prend $\mathcal{C} = \mathcal{C}'$, $u_! = u^* = \id_{\mathcal{C}}$, $W' = \operatorname{Is}(\mathcal{C})$, $W = \Fl(\mathcal{C})$, donc $u_W^*$ est le foncteur canonique
\[
\underbrace{\Pi\,\mathcal{C}}_{\text{groupoïde fondamental}} \qquad \longleftarrow \mathcal{C}\ .
\]
Supposons que $\mathcal{C}$ soit 1-connexe, i.e. $\Pi\,\mathcal{C}$ soit équivalente à $e$. Soit $e\in\Ob\mathcal{C}$ tel que $\operatorname{End}_{\mathcal{C}}(e)$ soit le groupe unité, et prenons $\mathcal{C}_0 = \{e\}$. Alors l'hypothèse sur $\mathcal{C}_0$ est satisfaite. Mais si $u_!^W$ était adjoint à gauche de $u_W^*$, $e$ serait objet initial de $\mathcal{C}$. Par exemple si $\mathcal{C}$ a un objet final $e$ (ce qui implique que $\mathcal{C}$ est 1-connexe et $e$ rigide), cet objet final devrait être initial si \textsc{Anderson} avait raison avec “it is clear that \ldots”.

\msp{page 60}
Néanmoins :

\textbf{Théorème 3.7.} \textit{Soit $(\M,W,C)$ une catégorie à cofibrations stable par $\varinjlim$ finies, et $u : I \longrightarrow I'$ une application croissante entre ensembles ordonnés artiniens. On suppose qu'on est sous l'une des hypothèses suivantes.}
\begin{enumerate}
\item[a)] $\Esp(I)$, $\Esp(I')$ \textit{sont noethériens, p.ex. $I$, $I'$ finis.}

\printedpage{30}
% head: A. Grothendieck : Dérivateurs, Ch. XVII, Édition des Universités de Montpellier II et Paris VII
% foot: Version 9 février 2019
\item[b)] \textit{$C$ et $TC \overset{\text{déf}}{=} C\cap W$ sont stables par composition transfinie, et $\M$ stable par petites $\varinjlim$ filtrantes.}
\end{enumerate}
\textit{Alors le foncteur} $u_!^W : \Ho_W(I^\circ) \longrightarrow \Ho_W(I'^\circ)$ \textit{construit dans} (3.3.1) \textit{est adjoint à gauche du foncteur $u_W^*$ en sens opposé.}

\textsc{Démonstration.} On utilise le diagramme bien connu
\[
\begin{tikzcd}[row sep=large, column sep=large]
 & I \arrow[d, "\partial_1"'] \arrow[dd, bend left=60, "w"] \\
I \arrow[r, shift left, "\partial_0\ (\text{ouv.})"] \arrow[d, "u"'] \arrow[dr, phantom, "\text{cocart.}"] & I\times\boldsymbol{\Delta}^1 \arrow[l, dotted, shift left, "\pr_1"] \arrow[d] \\
I' \arrow[r, shift left, "\text{imm. ouv.}"] & C(u) \arrow[l, dotted, shift left, "v"]
\end{tikzcd}
\]
pour factoriser $u$ en
\[
u \;=\; vw\ ,
\]
où $w$ est une \textit{immersion fermée,} et $v$ est déduit de $\pr_1 : I\times\boldsymbol{\Delta}^1 \longrightarrow I$ par push-out. C'est une rétraction du mapping-cone C(u) sur l'ouvert $I'$, dont le complémentaire fermé s'identifie à $I$ via $w$.

\msp{page 61}
Le foncteur de recollement entre $I'$, $I$
\[
H \;:\; I'^\circ\times I \;\longrightarrow\; \mathrm{Ens}
\]
est donné par
\[
\underbrace{H(i',i)}_{=\,\Hom_{C(u)}(i',i)} \;=\; \Hom_{I'}(i',u(i))\ . \tag*{$(*)$}
\]
La rétraction $v$ de $C(u)$ dans $I'$ applique $i$ dans $u(i)$, et une flèche $i'\longrightarrow i$ correspondant à $i'\longrightarrow u(i)$ dans $I'$ via $(*)$, est transformé par $v$ en cette même flèche $i'\longrightarrow u(i)$. Calculons, % [sic] transformé -> transformée, accord avec « une flèche »
pour $i'_0$ donné, la catégorie $i'_0\backslash\Bigl(C(u),\ v : C(u)\longrightarrow I'\Bigr)$, des couple d'un objet $\xi$ de $C(u)$ % [sic] des couple -> des couples, accord du pluriel
et d'une flèche $i'_0\longrightarrow v(\xi)$ dans $I'$. On voit que $i'_0\backslash C(u)$ admet l'ouvert $i'_0\backslash I'$, avec le complémentaire fermé $i'_0\backslash I$,
\[
\begin{tikzcd}[column sep=large, row sep=huge]
I' \arrow[r, hook, "\text{ouvert}"'] & C(u) \arrow[d] & I \arrow[l, hook', "\begin{array}{l}\text{fermé}\\ \text{complémentaire}\\ \text{de } I'\end{array}"'] \\
 & I'\ ,
\end{tikzcd}
\]
ces deux catégories étant recollées par le bifoncteur
\[
\begin{aligned}
(i'_0\backslash I')^\circ \times (i'_0\backslash I) &\quad\longrightarrow\quad \mathrm{Ens}\\
\Bigl((\alpha : i'_0\longrightarrow i'), (\beta : i'_0\longrightarrow u(i), i)\Bigr) &\quad\longmapsto\quad \Hom_{i'_0\backslash I'}(i',u(i))
\end{aligned}
\]
\printedpage{31}
% head: A. Grothendieck : Dérivateurs, Ch. XVII, Édition des Universités de Montpellier II et Paris VII
% foot: Version 9 février 2019
induit par $(*)$\,; i.e. les flèches de $(\alpha : i'_0\longrightarrow i') \in \Ob\,(i'_0\backslash I')$ dans $(\beta : i'_0\longrightarrow u(i)) \in i'_0\backslash I$ sont les flèches $i' \xrightarrow{\gamma} u(i)$ telles que le triangle
\[
\begin{tikzcd}
 & i'_0 \arrow[dl, "\alpha"'] \arrow[dr, "\beta"] & \\
i' \arrow[rr, "\gamma"] & & u(i)
\end{tikzcd}
\]
commute.

\msp{page 62}
Si $i' = i'_0$, $\alpha = \id_{i'_0}$, alors cette commutativité revient à $\gamma = \beta$, ce qui signifie qu'il y a une flèche et une seule de $\alpha$ dans $(i,\beta)$. Cela montre donc que $(i'_0,\id_{i'_0})$ est un objet initial de $i'_0\backslash C(u)$. On a donc prouvé le lemme :

\textbf{Lemme 3.8.} \textit{Soit $u : I\longrightarrow I'$ une flèche dans} (Cat)\textit{, factorisons-le en}
\[
I \xrightarrow{\;w\;} C(u) \xrightarrow{\;v\;} I'\ ,
\]
\textit{où $C(u)$ est le mapping-cone défini page} 60\textit{, $w$ une immersion fermée de complémentaire $I'$, $v$ la rétraction canonique de $C(u)$ sur $I'$. Ceci posé, pour tout $i'_0\in\Ob I'$, la catégorie $i'_0\backslash(C(u),v)$ admet un objet initial – donc $v : C(u)\longrightarrow I'$ est (tout ce qu'il y a de) cofiltrante.}

Nous pouvons maintenant prouver 3.7. Tout d'abord, les hypothèses faites sur $I$, $I'$ se transportent à $C(u)$ (hypothèse artinienne, et éventuellement hypothèse noethérienne sur $\Esp(C(u))$). Donc
\msp{page 63}
on peut appliquer les résultats précédents, notamment 3.5, aux flèches $v$, $w$. Le lemme précédent montre bien que $v$ satisfait aux conditions de 3.5, donc le foncteur $v_!^W$ construit dans 3.3.1 est adjoint à gauche de $v_W^*$. Considérons d'autre part le foncteur $w_!^W$ construit dans 3.3.1, je dis que la même conclusion s'applique. Cela provient du fait que $w$ est une immersion fermée. En effet, dans ce cas il est immédiat que
\[
w_! \;:\; \M(I^\circ) \;\longrightarrow\; \M(I'^\circ), \qquad\quad \text{où } I = C(u), % [sic] I = C(u) -> I' = C(u), w a pour but le mapping-cone C(u)
\]
est compatible avec les localiseurs, puisque l'on a
\[
\left\{
\begin{array}{lcll}
w_!(X)(j) & \simeq & X(j) & \text{si } j\in I,\\[1ex]
w_!(X)(j) & \simeq & \emptyset_{\M} & \text{si } j\in J\backslash I,
\end{array}
\right.
\]
par le calcul des fibres
\[
w_!(X)(j) \;\simeq\; \varinjlim_{i\in I/j} X(i)\ .
\]
Donc la définition de $w_!^W$ donnée dans 3.3.1 revient en fait à passer directement aux catégories localisées pour $w_!$, sans avoir à se restreindre à $\M(I^\circ)_{\mathrm{c}}$, et le lemme des foncteurs adjoints nous dit que
\printedpage{32}
% head: A. Grothendieck : Dérivateurs, Ch. XVII, Édition des Universités de Montpellier II et Paris VII
% foot: Version 9 février 2019
\msp{page 64}
ce foncteur est adjoint à gauche de $w_W^*$.

Comme on a évidemment
\[
u_!^W \;\simeq\; v_!^W\circ w_!^W\ , \qquad\quad u_W^* \;\simeq\; w_W^*\circ v_W^*\ ,
\]
on en conclut que $u_!^W$ est bien adjoint à gauche de $u_W^*$, q.e.d.

\textbf{Corollaire 3.9.} \textit{Soit $(\M,W,C)$ une catégorie à cofibrations. Alors le prédérivateur $\mathbf{D}_W$, restreint à la catégorie} $\mathrm{Diag}_0$ \textit{formée des catégories $I$} ordonnées \textit{artiniennes et telles que} $\Esp(I)$ \textit{soit noethérien, satisfait aux axiomes} Dér 1 (décomposition), Dér 2 (localisation), Dér 3 (objet final), Dér 4 bis (existence de $u_!$), Dér 5 bis (mode de calcul des $u_!$), Dér 6 (exactitude des carrés $W$-cocartésiens). \textit{Si $C$, $TC = C\cap W$ sont stables par composition transfinie, on a le même résultat sur la catégorie} Diag \textit{formée des catégories $I$ ordonnées artiniennes sans plus.}

\msp{page 65}
\textbf{3.10.} Ainsi, pour avoir un dérivateur au plein sens du terme, soit sur $\mathrm{Diag}_0$, soit sur Diag, il ne manque que Dér 4 (existence des $u_*$), Dér 5 (calcul des $u_*$) et Dér 6 (exactitude des carrés $W$-cartésiens). On va les obtenir en principe à l'aide d'une structure de \textit{catégorie à fibrations} $(\M,F,W)$ associé à $(\M,W)$, i.e. l'espèce de structure duale de celle intro- % [sic] associé -> associée, accord avec « structure »
duite page 39. Ces deux structures n'ont pas à être reliées entre elles d'une façon autre que par le fait qu'elles correspondent au même $W$. Elles sont considérablement plus faciles à construire que des structures de Quillen $(W,C,F)$, où $C$, $F$ doivent être reliés de façon extrêmement stricte, avec les propriétés $C \longleftrightarrow \underbrace{TF}_{=\,F\cap W}\ ,\ \underbrace{TC}_{=\,C\cap W} \longleftrightarrow F$ (dont il n'est pas question dans l'axiomatique de K. S. \textsc{Brown}), et factorisation pour $(C,TF)$ et $(TC,F)$.

\msp{page 66}
Mais revenant à la construction des $u_!^W$, il reste le problème de les construire pour une flèche $u : I'\longrightarrow I$ \textit{quelconque} dans Cat. % [sic] u : I' -> I -> u : I -> I', u_!^W va de Ho_W(I°) dans Ho_W(I'°)

À vrai dire, la construction du foncteur
\[
u_!^W \;:\; \Ho_W(I^\circ) \;\longrightarrow\; \Ho_W(I'^\circ) \tag*{(3.3.1)}
\]
est valable sans supposer que $I'$ soit une catégorie ordonnée – ça peut être une catégorie quelconque. Cela provient du fait que le foncteur induit par $u_!$,
\[
u_{!\,\mathrm{cof}} \;:\; \M(I^\circ)_{\mathrm{c}} \;\longrightarrow\; \M(I'^\circ)
\]
est encore compatible avec les localiseurs, comme on voit par le calcul des fibres
\[
u_!(X)(i) \;\simeq\; \varinjlim_{i\in I/i'} X(i)\ . % [sic] u_!(X)(i) -> u_!(X)(i'), la fibre est prise en i'
\]
Ainsi, on est ramené aux cas où $I' = e$, quitte à remplacer $I$ par $I/i'$, et en notant que % [sic] aux cas -> au cas, singulier
$(I/i'\longrightarrow I$ étant étale) l'image inverse de $X$ sur $I/i'$ est cofibrant (cf. 3.2.a). % [sic] cofibrant -> cofibrante, accord avec « l'image inverse »

\msp{page 67}
Mais il se présente à nouveau la perplexité si le foncteur (3.3.1) ainsi construit est bien adjoint à gauche de $u_W^*$. On peut songer à utiliser encore la même factorisation en $u = vw$, en
\printedpage{33}
% head: A. Grothendieck : Dérivateurs, Ch. XVII, Édition des Universités de Montpellier II et Paris VII
% foot: Version 9 février 2019
utilisant le mapping-cone. La difficulté, cependant, c'est qu'à présent $J = C(u)$ n'est plus une catégorie ordonnée, donc on n'a pas une définition toute prête de $C_{J^\circ}\subset\Fl(\M(J^\circ))$, donnant lieu à une sous-catégorie $\M(J^\circ)_{\mathrm{c}}$ etc. En fait, il semblerait qu'on soit pratiquement paralysé, tant qu'on n'arrive à construire sur $\M(J^\circ)$ également une structure de catégorie à cofibrations, se comportant bien par rapport au foncteurs $v_!$. La % [sic] au foncteurs -> aux foncteurs, accord du pluriel
“démonstration” d'\textsc{Anderson} de l'existence de $u_!$ pour \textit{toute} flèche dans Cat, foire totalement (même pour $I\longrightarrow e$, $I$ ensemble ordonné fini – mais là j'ai pu rétablir la situation), car il se ramène au cas où $I$ est ordonné localement fini, et il
\msp{page 68}
admet sans plus que \textit{dans ce cas} $u_!^W$ est adjoint à gauche de $u_W^*$, et \textit{de plus,} que si $u$ est $\mathbf{W}_0$-connexe, donc $u^* : \M(I'^\circ)\longrightarrow\M(I^\circ)$ pleinement fidèle, qu'il en soit de même de $u_W^*$ déduit par passage aux localisés. Or même si on admet que $u_!^W$ est adjoint à gauche de
$u_*^W$, faute de disposer d'un morphisme d'adjonction tout près “en bas” (dans $\Ho_W(I'^\circ)$), % [sic] u_*^W -> u_W^*, notation employée partout ailleurs
déduit du morphisme d'adjonction dans $\M(I^\circ)$, il n'est absolument pas clair que $u_W^*$ soit pleinement fidèle, i.e. $u_!^W$ un foncteur de localisation. Le contrexemple 3.6 montre bien que, même si on a adjonction (i.e. si $e$ est un objet initial de $\mathcal{C} = \mathcal{C}'$), le foncteur $u^* = \id_{\mathcal{C}}$ a beau être pleinement fidèle (et même une équivalence de catégories), le foncteur localisé n'a aucune raison de l'être : cela pouverait que tout objet de Cat ayant un objet initial % [sic] pouverait -> prouverait, faute de frappe
est équivalent à $e$\,!

\msp{page 69}
\textbf{3.11.} Disons qu'on a une \textit{catégorie à cofibrations stricte,} si dans l'axiome $\mathrm{Cof}_3$ de factorisation (cf. p. 39), la factorisation peut être prise fonctoriellement en $f$. Il revient $\pm$ au même de dire que l'axiome de factorisation reste valable, quand on pense à une catégorie $\M(I)$ pour $I$ dans Cat, en y définissant $W(I)$ et $C(I)$ (à ne pas confondre avec $C_I$) fibre par fibre. Cela revient au même tout au moins si $\M$ est équivalente à une petite catégorie. Sinon, il faut supposer fonctorialité pour toute \textit{petite} sous-catégorie pleine de $\underline{\mathrm{Fl}}(\M)$. Il vaut mieux définir la catégorie à cofibrations stricte par cette propriété de stabilité.

On voit alors que le prédérivateur $\mathbf{D} = \mathbf{D}_{\M}$ sur Cat est un dérivateur \textit{relativement à} $\mathcal{D} = \mathrm{Diag}_0$ défini plus haut (catégories ordonnées $I$ artiniennes, telles que $\Esp(I)$ soit noethérien), ou même relativement à $\mathcal{D} = \mathrm{Diag}$ (catégories ordonnées $I$ artiniennes sans plus) lorsque $C$ et $TC$ sont stables par composition transfinie. J'entends par là que $u_!$ existe quand $u$ est une flèche de $\mathcal{D}$,
\msp{page 70}
et que ça reste valable quand on passe à un dérivateur induit $\mathbf{D}_Z$ ($Z$ dans Cat quelconque)\,; en d'autre termes, que $v_!$ existe quand $v$ est de la forme $\id_Z\times u$, où $u : I\longrightarrow I'$ est une % [sic] d'autre termes -> d'autres termes, accord du pluriel
flèche de $\mathcal{D}$, et que le calcul se fait \textit{au dessus de} $I'$ selon la règle Dér 5 bis. Cela signifie donc que la flèche de changement de base pour les $f_!$ est un isomorphismes dans le cas % [sic] un isomorphismes -> un isomorphisme, accord du singulier
suivant :
\[
\begin{tikzcd}[row sep=huge, column sep=huge]
I\times Z \arrow[d, "u\times\id_Z"] & (i'\backslash I)\times Z \arrow[l] \arrow[d, "\pr_2"] \\
I'\times Z & \{i'\}\times Z \arrow[l]
\end{tikzcd}
\ ,
\]
où $i'$ est un objet quelconque de $I'$.

\end{document}
